\chapter{Introduction}

Software engineering has always been closely tied to physics and mathematics to the point where the former has become strongly associated with 
the latter. In some fields, such as game development which has a very large indepandant / solo community, this may be considered to be a barrier to 
entry for some. People may be dissuaded from learning how to program due to how heavy the mathematics might be, leading to the field missing out on 
great games made from the passion of talented people.

Resources such as research papers which are not explicitly designed with beginners-to-the-field, or game developers mystify the mathematics due to
heavily verbose notation and content.

\section{Project Aim}

This project aims to evaluate and review how approachable research papers are for people who do not have a background in mathematics and physics,
in creating a project.

The aim of this project is to create a plausible simulation of particle physics in C++, a project that will be heavily based in foundational physics.

\section{Subjectivity Disclaimer}

The evaluation of the accessibility and approachability of scientific papers and by extension the developer experience detailed
in this report is inherently subjective, shaped by my own specific background in both particle physics (or lack thereof),
and experience in programming, and interpretation of the resources used. As such, the challenges and successes documentated
will vary for other developers with different backgrounds and levels of expertise in the two fields.